%% kapa: exam=1 round=23 year=2012 period=1 form=A subject=civil_law subject_ko=민법 seq=23 number=23
%% source: Q-Net 감정평가사 기출문제 (cst00309) artlSeq=5052578 file=raw/1008/5052578/2012년 제23회 감정평가사 제1차시험 기출문제(1교시 A형).pdf
%% conversion: kordoc 4.24.51 -> tools/gen1.py (auto); review=proofread-sonnet; flags=none
\begin{question}{23}
乙은 대리권이 없음에도 불구하고 甲의 대리인이라고 하면서 甲 소유의 X 토지를 丙에게 매도하는 계약을 체결하였다. 다음 설명 중 옳은 것은?
\begin{choices}
  \item 甲이 乙에 대하여 추인한 사실을 丙이 알았다고 하더라도 甲은 丙에 대하여 추인의 효과를 주장할 수 없다.
  \item 丙이 乙의 대리권 없음을 알고 있었다면, 丙은 甲에 대하여 추인 여부의 확답을 최고할 수 없다.
  \item 丙은 乙의 대리권 없음을 알지 못한 데 대하여 과실이 있는 경우에도 乙에게 계약의 이행을 청구할 수 있다.
  \item 丙이 乙의 대리권 없음을 알지 못한 경우에는 乙의 선택에 따라 계약의 이행을 청구할 수 있다.
  \item 丙이 乙의 대리권 없음을 알지 못한 경우에는 계약을 철회할 수 있다.
\end{choices}
\end{question}
\answer{5}
